%% CAPITULO 5 -- CONSIDERACOES FINAIS (esqueleto da dissertação)

\section{Considerações finais}

Espera-se que a principal contribuição desta dissertação seja um instrumento
validado para medir a rastreabilidade de respostas geradas por sistemas RAG
sobre documentos públicos, complementando métricas técnicas que avaliam
relevância e fidelidade, mas não a indicação verificável da fonte. A matriz de
decisão, como contribuição secundária, oferece ao comitê de governança de dados
e de inteligência artificial um critério explícito para escolher entre
estratégias de contextualização.

O estudo tem limites assumidos desde o projeto: o corpus é brasileiro e reúne um
conjunto delimitado de tipos documentais, o modelo gerador é mantido fixo e a
avaliação do artefato não inclui o uso real em um órgão. A avaliação em uso, a
transferência do IRD para outros tipos documentais e a análise de equidade entre
públicos ficam como desdobramentos para pesquisas futuras.
